\chapter{Четырёхмерная суперсимметрия и модель Wess--Zumino}

\lectureoverview{Одномерная алгебра суперсимметрии из предыдущей лекции включала только гамильтониан. Теперь инвариантность Лоренца требует снова объединить энергию с пространственным импульсом, а суперзаряды должны стать двухкомпонентными спинорами. К этой алгебре мы придём, сначала вернувшись к безмассовому фермиону Вейля и его гамильтониану \(H=\boldsymbol{\sigma}\cdot\mathbf p\). Затем мы представим суперзаряды как дифференциальные операторы в четырёхмерном суперпространстве, выделим киральные суперполя с помощью ограничения \(\bar D_{\dot\alpha}\Phi=0\) и воспользуемся интегрированием по суперпространству, чтобы построить конкретную теорию поля. Последовательный выбор \(\Phi^\dagger\Phi\), \(\Phi^2\) и \(\Phi^3\) обеспечивает распространение, равные бозон--фермионные массы, юкавские взаимодействия и связанное с ними скалярное взаимодействие. Вспомогательное поле \(F\) делает эти соотношения наглядными, а заключительные петлевые диаграммы показывают, почему суперсимметрия смягчает чувствительность в ультрафиолете.}

\section{Почему гамильтониан не может стоять сам по себе}

В игрушечной модели пространство было опущено, и основное соотношение имело вид \begin{equation}
\{Q,Q^\dagger\}=2H.
\label{eq:l8-time-only-algebra}
\end{equation} Эта форма достаточна, когда время является единственной координатой. Однако в релятивистской теории \(H=P_0\) — лишь одна из компонент четырёхимпульса \(P_\mu\). Буст Лоренца смешивает её с \(P_i\). Алгебра, правая часть которой выделяла бы \(P_0\), не могла бы преобразовываться ковариантно.

Следовательно, исправление должно одновременно сделать две вещи. Оно должно заменить \(H\) всеми четырьмя компонентами \(P_\mu\) и должно обеспечить слева индексы, которые можно сопоставить с лоренцевым вектором справа. Нужную подсказку даёт безмассовое уравнение Дирака: релятивистский фермион можно описать двухкомпонентным спинором, а матрицы \(\sigma^\mu\) связывают один векторный индекс с парой спинорных индексов.

\section{Безмассовые фермионы, спиральность и киральность}

Начнём с физического различия между массивными и безмассовыми частицами. Рассмотрим массивную частицу со спином \(\tfrac12\), движущуюся в направлении своего спина. Поскольку её скорость меньше скорости света, наблюдатель может её обогнать. В новой системе отсчёта импульс меняет знак, а спин — нет, так что спиральность меняется на противоположную. Поэтому лоренц-инвариантная теория массивной частицы должна содержать обе спиральности.

Никакой наблюдатель не может обогнать безмассовую частицу. Собственные преобразования Лоренца не могут обратить её импульс относительно спина, поэтому одно неприводимое безмассовое представление может иметь одну спиральность. Пространственный поворот может развернуть направление движения, но не превращает правую руку в левую.

\begin{classroomqa}\audiencequestion{Действительно ли аргумент с обгоном относится к спиральности, а не к киральности? \lecturetimestamp{00:05:42}}\lecturerresponse{Да. Спиральность — это проекция спина на импульс, тогда как киральность — собственное значение киральной матрицы, обычно \(\gamma^5\). Для безмассовой моды эти две метки связаны друг с другом, и именно поэтому в этом примере язык может звучать так, словно они взаимозаменяемы. Для массивной частицы они различны.}\end{classroomqa}

Число состояний требует одной оговорки.\editorialnote{Одно безмассовое поле Вейля имеет одну спиральность у частицы и противоположную — у античастицы. Поэтому «вдвое меньше состояний» относится к переходу от четырёхкомпонентного поля Дирака к одному двухкомпонентному полю Вейля; это не означает, что состояния античастицы исчезли.}

Поле Вейля, используемое в этом начальном аргументе, запишем как \begin{equation}
\chi=
\begin{pmatrix}
\chi_1\\
\chi_2
\end{pmatrix},
\end{equation} В лекции это поле обозначается как \(\psi_\alpha\); здесь мы будем называть его \(\chi\), чтобы не путать с сопряжённым представлением Вейля, которое по соглашению ниже помещается в киральное суперполе. \begin{equation}
H\chi=\boldsymbol{\sigma}\cdot\mathbf p\,\chi.
\label{eq:l8-weyl-hamiltonian}
\end{equation} Для импульса \(\mathbf p\) матрица \(\boldsymbol{\sigma}\cdot\mathbf p\) имеет собственные значения \(\pm|\mathbf p|\). Положительно-частотная ветвь даёт моды частицы; отрицательно-частотная ветвь переинтерпретируется через построение античастицы, а не сохраняется как спектр, не ограниченный снизу.

\begin{figure}[H]\centering
\includegraphics[width=0.70\textwidth]{lecture_08_figure_02.png}\lectureframe{00:07:44.271}\caption{Киральный двухкомпонентный спинор и гамильтониан \(H=\boldsymbol{\sigma}\cdot\mathbf p\). Последний символ импульса частично закрыт, но произнесённое уравнение его фиксирует.}\end{figure}

Обозначение рукости также зависит от соглашения.\editorialnote{На доске рукость напрямую связывается с положительным собственным значением \(\boldsymbol{\sigma}\cdot\mathbf p\). То, какая киральность получает этот знак, зависит от определений \(\sigma^\mu\), \(\bar\sigma^\mu\) и поля Вейля. Мы сохраняем утверждение о собственном значении, но не приписываем ему не зависящую от соглашения метку левой или правой рукости.}

Замена \(H\) на \(i\partial_t\) и \(\mathbf p\) на \(-i\nabla\) даёт \begin{equation}
i\partial_t\chi=-i\boldsymbol{\sigma}\cdot\nabla\chi.
\end{equation} Введём \begin{equation}
\begin{aligned}
\sigma^\mu&=(\mathbf 1,\boldsymbol{\sigma}),\\
\bar\sigma^\mu&=(\mathbf 1,-\boldsymbol{\sigma}),\\
\eta_{\mu\nu}&=\operatorname{diag}(1,-1,-1,-1).
\end{aligned}
\label{eq:l8-sigma-convention}
\end{equation} Тогда только что полученное уравнение принимает вид \begin{equation}
i\sigma^\mu\partial_\mu\chi=0.
\label{eq:l8-weyl-equation}
\end{equation} Сопряжённое представление Вейля использует \(\bar\sigma^\mu\). Именно это сопряжённое представление, обозначаемое ниже как \(\psi_\alpha\), имеет свободное уравнение \(i\bar\sigma^\mu\partial_\mu\psi=0\). Доказательство ковариантности требует спинорного представления группы Лоренца, но здесь важен структурный факт: четыре матрицы \(\sigma^\mu\) связывают один пространственно-временной индекс с двумя спинорными индексами.

\section{Четырёхмерная супералгебра}

Теперь нечётные координаты должны преобразовываться как спиноры Вейля: \begin{equation}
\theta^\alpha=(\theta^1,\theta^2),
\qquad
\bar\theta^{\dot\alpha}=(\bar\theta^{\dot1},\bar\theta^{\dot2}).
\end{equation} Черта обозначает комплексное сопряжение, а точка отличает сопряжённое представление Лоренца. Имеются две комплексные, то есть четыре вещественные, грассмановы компоненты. Поэтому общее чётное суперполе \begin{equation}
\mathcal S=\mathcal S(x^\mu,\theta^\alpha,\bar\theta^{\dot\alpha})
\end{equation} имеет конечное разложение по переменным Грассмана. Его моном наивысшей степени имеет вид \begin{equation}
\theta^2\bar\theta^2
\equiv
(\theta^1\theta^2)(\bar\theta^{\dot1}\bar\theta^{\dot2}),
\end{equation} Это произведение всех четырёх независимых нечётных координат. Коэффициент при этом мономе обычно называют \(D\)-компонентой.

Суперзаряды тоже несут спинорные индексы. В четырёхмерной суперсимметрии \(N=1\) их алгебра имеет вид \begin{align}
\{Q_\alpha,\bar Q_{\dot\beta}\}
&=2\sigma^\mu_{\alpha\dot\beta}P_\mu,
\label{eq:l8-superalgebra-main}\\
\{Q_\alpha,Q_\beta\}&=0,
&
\{\bar Q_{\dot\alpha},\bar Q_{\dot\beta}\}&=0,
\label{eq:l8-superalgebra-nilpotent}\\
[P_\mu,Q_\alpha]&=0,
&
[P_\mu,\bar Q_{\dot\alpha}]&=0,
&
[P_\mu,P_\nu]&=0.
\label{eq:l8-superalgebra-commutators}
\end{align} Множитель 2 определяется соглашением. Матрица \(\sigma^\mu_{\alpha\dot\beta}\) свёртывает векторный индекс у \(P_\mu\) и оставляет ровно те два спинорных индекса, которые стоят слева.

\begin{classroomqa}\audiencequestion{Являются ли \(\alpha\) и \(\beta\) у сигма-матрицы её матричными индексами? \lecturetimestamp{00:22:01}}\lecturerresponse{Да. Точнее, один из них — неточечный индекс Вейля, а другой — точечный индекс сопряжённого представления; \(\mu\) — это пространственно-временной векторный индекс. Сигма-матрица служит мостом между этими описаниями.}\end{classroomqa}

Отредактированная нотация делает этот мост явным.\editorialnote{В лекции точечные индексы намеренно опускаются, а на доске обсуждается порядок \(\alpha\beta\) по сравнению с \(\beta\alpha\). Уравнения~\eqref{eq:l8-superalgebra-main} и \eqref{eq:l8-superalgebra-commutators} восстанавливают стандартную нотацию с точечными и неточечными индексами, не меняя приводимого аргумента.}

Последняя строка означает, что суперзаряды сохраняются и совместимы с трансляциями. Компоненты импульса коммутируют, потому что коммутируют их дифференциальные представители. В дифференциальном соглашении, используемом ниже, \begin{equation}
P_\mu=i\partial_\mu,
\end{equation} поэтому соотношение~\eqref{eq:l8-superalgebra-main} с тем же успехом можно записать как \begin{equation}
\{Q_\alpha,\bar Q_{\dot\beta}\}
=2i\sigma^\mu_{\alpha\dot\beta}\partial_\mu.
\end{equation}

\section{Суперзаряды как производные в суперпространстве}

Импульс порождает трансляцию обычной координаты; суперзаряд порождает трансляцию нечётной координаты, одновременно сдвигая пространство-время. Выберем левые производные по переменным Грассмана и определим \begin{align}
Q_\alpha
&=\partial_\alpha
-i\sigma^\mu_{\alpha\dot\alpha}
\bar\theta^{\dot\alpha}\partial_\mu,
\label{eq:l8-Q}\\
\bar Q_{\dot\alpha}
&=-\bar\partial_{\dot\alpha}
+i\theta^\alpha\sigma^\mu_{\alpha\dot\alpha}\partial_\mu.
\label{eq:l8-Qbar}
\end{align} Действие \(\epsilon^\alpha Q_\alpha+
\bar\epsilon_{\dot\alpha}\bar Q^{\dot\alpha}\) даёт бесконечно малую вариацию суперсимметрии суперполя. Его компоненты — это просто коэффициенты при независимых мономах Грассмана, а дифференциальные операторы перемешивают эти коэффициенты.

\begin{classroomqa}\audiencequestion{Какая производная по \(x\) имеется в виду в сокращённой записи \(\bar\theta\sigma\partial\)? \lecturetimestamp{00:25:40}}\lecturerresponse{Все четыре. Повторяющийся индекс \(\mu\) в \(\bar\theta^{\dot\alpha}\sigma^\mu_{\alpha\dot\alpha}\partial_\mu\) суммируется, так что каждая компонента \(\sigma^\mu\) спаривается с соответствующей производной \(\partial/\partial x^\mu\).}\end{classroomqa}

\begin{classroomqa}\audiencequestion{Не должен ли сопряжённый заряд содержать \(\theta\), а не ещё один \(\bar\theta\)? \lecturetimestamp{00:26:16}}\lecturerresponse{Да. \(Q\) дифференцирует по \(\theta\) и содержит \(\bar\theta\); \(\bar Q\) дифференцирует по \(\bar\theta\) и содержит \(\theta\). Это сопряжённые операторы.}\end{classroomqa}

\begin{classroomqa}\audiencequestion{Является ли \(\sigma^\mu\) симметричной? \lecturetimestamp{00:27:04}}\lecturerresponse{Матрицы Паули эрмитовы. Три из четырёх матриц в \((\mathbf1,\boldsymbol{\sigma})\) симметричны, тогда как \(\sigma^2\) антисимметрична при обычном транспонировании. Эрмитовость вместе с расстановкой точечных и неточечных индексов и есть здесь существенное свойство сопряжения.}\end{classroomqa}

\begin{classroomqa}\audiencequestion{Несёт ли сигма-матрица в сопряжённом уравнении индексы в обратном порядке? \lecturetimestamp{00:27:13}}\lecturerresponse{Сопряжение обменивает точечные и неточечные представления, так что безопасная нотация — \(\sigma^\mu_{\alpha\dot\alpha}\) и её сопряжение \(\bar\sigma^{\mu\,\dot\alpha\alpha}\), а не трактовка этих двух меток как взаимозаменяемых матричных позиций.}\end{classroomqa}

Формулы в уравнениях~\eqref{eq:l8-Q} и \eqref{eq:l8-Qbar} отличаются лишь внешними знаками от сокращённой записи, использованной на доске. Эти знаки связаны с тем, является ли \(P_\mu\) величиной \(+i\partial_\mu\) или \(-i\partial_\mu\), действуют ли производные слева или справа и в каком порядке записываются сопряжённые нечётные переменные. Важно то, что одно выбранное соглашение удовлетворяет всей алгебре сразу.

\section{Ограничение, сохраняющееся при суперсимметрии}

Общее суперполе содержит много компонент. Мы ищем меньший класс, не разрушая симметрию. Сначала в лекции приводится обычная аналогия. Для вектора \(\mathbf V\) ограничение
\begin{equation}
V_x=0
\end{equation}
не инвариантно относительно вращений: после поворота \(V_y\) даёт вклад в новую \(x\)-компоненту. Напротив,
\begin{equation}
\mathbf V^2=1
\end{equation}
инвариантно и устраняет одну независимую степень свободы. Полезное ограничение должно при действии симметрии переводить поля, удовлетворяющие ограничению, в поля, которые ему по-прежнему удовлетворяют.

Пусть линейный нечётный оператор \(\bar D_{\dot\alpha}\) задаёт
\begin{equation}
\bar D_{\dot\alpha}\Phi=0.
\label{eq:l8-chiral-constraint}
\end{equation}
После вариации суперсимметрии \(\delta\Phi=Q\Phi\) то же условие сохранится, если \(\bar D\) антикоммутирует с суперзарядами:
\begin{equation}
\bar D(Q\Phi)=-Q(\bar D\Phi)=0.
\end{equation}
Именно эта абстрактная проверка согласованности мотивирует построение, а не определение, появляющееся без объяснения.

\begin{classroomqa}
\audiencequestion{Почему в последнем выражении можно переставить \(Q\) и \(\bar D\)? \lecturetimestamp{00:36:56}}
\lecturerresponse{Это нечётные операторы, выбранные так, чтобы антикоммутировать. Поэтому перестановка их порядка вносит знак минус; как только \(\bar D\) доходит до \(\Phi\), исходное ограничение обращает результат в нуль.}
\end{classroomqa}

Требуемые ковариантные производные получаются обращением знаков при пространственно-временных производных по сравнению с суперзарядами:
\begin{align}
D_\alpha
&=\partial_\alpha
+i\sigma^\mu_{\alpha\dot\alpha}
\bar\theta^{\dot\alpha}\partial_\mu,
\label{eq:l8-D}\\
\bar D_{\dot\alpha}
&=-\bar\partial_{\dot\alpha}
-i\theta^\alpha\sigma^\mu_{\alpha\dot\alpha}\partial_\mu.
\label{eq:l8-Dbar}
\end{align}
Они удовлетворяют
\begin{align}
\{D_\alpha,\bar D_{\dot\beta}\}
&=-2i\sigma^\mu_{\alpha\dot\beta}\partial_\mu,\\
\{D_\alpha,Q_\beta\}
&=\{D_\alpha,\bar Q_{\dot\beta}\}=0,\\
\{\bar D_{\dot\alpha},Q_\beta\}
&=\{\bar D_{\dot\alpha},\bar Q_{\dot\beta}\}=0.
\end{align}

\begin{classroomqa}
\audiencequestion{Верно ли, что при переходе от \(\bar Q\) к \(\bar D\) меняется только знак у члена с пространственно-временной производной? \lecturetimestamp{00:37:58}}
\lecturerresponse{Именно. Это изменение знака приводит к сокращению перекрёстных членов, когда \(D\) антикоммутирует с любым из суперзарядов.}
\end{classroomqa}

Мы налагаем только условие ~\eqref{eq:l8-chiral-constraint}. Наложение сразу двух условий \(D_\alpha\Phi=0\) и \(\bar D_{\dot\alpha}\Phi=0\) через их антикоммутатор повлекло бы \(\partial_\mu\Phi=0\), оставляя самое большее остаток, не зависящий от пространства-времени.
\editorialnote{На доске сказано, что наложение обоих ограничений влечёт \(\Phi=0\). Строго говоря, алгебра требует лишь независимости от пространства-времени; дополнительные граничные или компонентные условия затем могут ещё сильнее ограничить поле.}

\section{Решение кирального ограничения}

Определим сдвинутую координату
\begin{equation}
y^\mu=x^\mu+i\theta\sigma^\mu\bar\theta.
\label{eq:l8-chiral-coordinate}
\end{equation}
При принятой выше конвенции левых производных функция
\begin{equation}
\Phi=\Phi(y,\theta)
\label{eq:l8-chiral-solution}
\end{equation}
удовлетворяет \(\bar D_{\dot\alpha}\Phi=0\). Правило цепочки делает сокращение явным. Дифференцирование через \(y^\mu\) даёт тот же член \(i\theta\sigma^\mu\partial_\mu\Phi\), который входит в \(\bar D_{\dot\alpha}\) с противоположным знаком.

Таким образом, киральное суперполе не имеет независимой зависимости от \(\bar\theta\). Оно всё же может содержать \(\bar\theta\), но только через комбинацию \(y^\mu\). Его конечное разложение имеет вид
\begin{equation}
\Phi(y,\theta)
=\phi(y)+\sqrt2\,\theta\psi(y)+\theta^2F(y).
\label{eq:l8-chiral-expansion-y}
\end{equation}
Три типа компонент — это комплексный скаляр \(\phi\), вейлевский фермион \(\psi_\alpha\) и комплексный скаляр \(F\). Последний является вспомогательным, а не вторым распространяющимся бозоном.

\begin{classroomqa}
\audiencequestion{Была ли \(x\) заменена новой координатой \(y\), ортогональной ей? \lecturetimestamp{00:49:08}}
\lecturerresponse{Нет. \(y^\mu\) — это сокращённое обозначение конкретной комбинации из ~\eqref{eq:l8-chiral-coordinate}; это не дополнительное пространственно-временное направление. Работа в \(y\) делает киральность простой, а при построении обычной теории поля мы возвращаемся к \(x\).}
\end{classroomqa}

Разлагая уравнение ~\eqref{eq:l8-chiral-expansion-y} при фиксированном \(x\), получаем
\begin{align}
\Phi(x,\theta,\bar\theta)
={}&\phi
+\sqrt2\,\theta\psi
+\theta^2F
+i\theta\sigma^\mu\bar\theta\,\partial_\mu\phi
\notag\\
&-\frac{i}{\sqrt2}\theta^2
\partial_\mu\psi\,\sigma^\mu\bar\theta
+\frac14\theta^2\bar\theta^2\Box\phi.
\label{eq:l8-chiral-expansion-x}
\end{align}
Все компонентные поля в правой части берутся в точке \(x\). Сопряжённое антихиральное поле зависит от
\(\bar y^\mu=x^\mu-i\theta\sigma^\mu\bar\theta\) и \(\bar\theta\). Киральное суперполе не может согласованно быть вещественным, поскольку его сопряжённое удовлетворяет условию противоположной киральности.

\section{Интегралы по суперпространству}

Действие по полному суперпространству имеет вид
\begin{equation}
S_D=\int d^4x\,d^2\theta\,d^2\bar\theta\;K,
\label{eq:l8-D-term-action}
\end{equation}
где \(K\) — вещественное суперполе. Интегрирование Березина выделяет коэффициент при \(\theta^2\bar\theta^2\). После интегрирования по нечётным координатам остаётся только обычный лагранжиан пространства-времени.

Существует и вторая инвариантная конструкция. Если \(W(\Phi)\) кирально, то
\begin{equation}
S_F=\int d^4x\,d^2\theta\;W(\Phi)
+\int d^4x\,d^2\bar\theta\;\overline{W}(\Phi^\dagger)
\label{eq:l8-F-term-action}
\end{equation}
является суперсимметричным. Первый интеграл выделяет компоненту \(\theta^2\), а второй — её сопряжённую.

\begin{classroomqa}
\audiencequestion{Почему допустимо интегрировать только по \(d^2\theta\), если выражение не зависит от \(\bar\theta\)? \lecturetimestamp{01:16:30}}
\lecturerresponse{Киральность сохраняется суперсимметрией, а старшая \(\theta\)-компонента кирального суперполя изменяется лишь на полную пространственно-временную производную. Поэтому её интеграл по пространству-времени инвариантен, точно так же как инвариантен полный интеграл по \(D\)-компоненте.}
\end{classroomqa}

\section{Свободный киральный мультиплет}

Рассмотрим выражение, на первый взгляд лишённое структуры,
\begin{equation}
S_{\mathrm{free}}
=\int d^4x\,d^2\theta\,d^2\bar\theta\;
\Phi^\dagger\Phi.
\label{eq:l8-free-superaction}
\end{equation}
Эти два множителя зависят от координат, сдвинутых в противоположные стороны:
\begin{equation}
\Phi=\Phi(x+i\theta\sigma\bar\theta,\theta),
\qquad
\Phi^\dagger
=\Phi^\dagger(x-i\theta\sigma\bar\theta,\bar\theta).
\end{equation}
Под знаком пространственно-временного интеграла мы можем сдвинуть фиктивную переменную \(x\). Один сдвиг может убрать грассманово смещение из одного множителя, но тогда он удваивает относительное смещение в другом. Разложение Тейлора по этому нильпотентному смещению обрывается после конечного числа членов.

Именно отсюда берутся производные. Суперпространственное выражение \(\Phi^\dagger\Phi\) не содержит явной пространственно-временной производной, однако сдвинутые аргументы порождают первые и вторые производные компонентных полей. Перемножая разложения и оставляя только старшую компоненту, получаем, в соглашении ~\eqref{eq:l8-sigma-convention},
\begin{equation}
\mathcal L_{\mathrm{free}}
=\partial_\mu\phi^\ast\partial^\mu\phi
+i\bar\psi\bar\sigma^\mu\partial_\mu\psi
+F^\ast F.
\label{eq:l8-free-component-lagrangian}
\end{equation}
С точностью до интегрирования по частям скалярную часть можно также схематически записать как \(-\phi^\ast\Box\phi\); именно к этой форме подходит вычисление на доске.

\begin{classroomqa}
\audiencequestion{Откуда берётся вторая строка в тейлоровском разложении суперполя? \lecturetimestamp{01:03:00}}
\lecturerresponse{Она возникает из компонентного разложения \(\Phi(y,\theta)=\phi(y)+\sqrt2\theta\psi(y)+\theta^2F(y)\). Затем каждую компоненту разлагают по Тейлору, поскольку \(y\) отличается от \(x\) нильпотентной грассмановой билинейной комбинацией.}
\end{classroomqa}

\begin{classroomqa}
\audiencequestion{Должна ли эта последняя компонента нести \(\theta^2\)? \lecturetimestamp{01:04:54}}
\lecturerresponse{Да. Вспомогательная компонента имеет вид \(F\theta^2\). Поскольку в ней уже содержатся обе независимые \(\theta\), любой дальнейший тейлоровский сдвиг, несущий ещё одну \(\theta\), обращается в нуль.}
\end{classroomqa}

\begin{figure}[H]
\centering
\includegraphics[width=0.82\textwidth]{lecture_08_figure_03.png}
\lectureframe{01:08:57.900}
\caption{Разложение по сдвинутой координате, собранное в скалярные, фермионные и вспомогательные члены обычного лагранжиана.}
\end{figure}

Первые два члена в уравнении ~\eqref{eq:l8-free-component-lagrangian} дают безмассовые уравнения Клейна--Гордона и Вейля. Последний не содержит производной:
\begin{equation}
\frac{\partial\mathcal L}{\partial F}=F^\ast=0,
\qquad
\frac{\partial\mathcal L}{\partial F^\ast}=F=0.
\label{eq:l8-free-F-equation}
\end{equation}
Он не задаёт ни начальных данных, ни частицы на массовой оболочке.

\begin{figure}[H]
\centering
\includegraphics[width=0.82\textwidth]{lecture_08_figure_04.png}
\lectureframe{01:11:32.360}
\caption{Вспомогательный член без производных \(F^\ast F\) в конце компонентного лагранжиана.}
\end{figure}

Язык пропагаторов говорит то же самое. Обратный квадратичный оператор для \(F\) не зависит от импульса. Его преобразование Фурье локально:
\begin{equation}
\langle F(x)F^\ast(y)\rangle
\propto \delta^{(4)}(x-y).
\label{eq:l8-F-contact}
\end{equation}
Линия \(F\) может служить для организации диаграммы, но её концы совпадают. Это контактная контракция, а не распространение физической частицы.

\section{Квадратичный суперпотенциал и равные массы}

Распространение возникало из \(D\)-члена. Масса возникает из кирального интеграла. Выберем стандартную нормировку
\begin{equation}
W_2(\Phi)=\frac12m\Phi^2.
\end{equation}
\(\theta^2\)-компонента равна
\begin{equation}
\left.W_2(\Phi)\right|_{\theta^2}
=m\phi F-\frac12m\psi\psi.
\end{equation}
Вместе с сопряжённым членом компонентный лагранжиан принимает вид
\begin{align}
\mathcal L
={}&
\partial_\mu\phi^\ast\partial^\mu\phi
+i\bar\psi\bar\sigma^\mu\partial_\mu\psi
+F^\ast F
\notag\\
&+m\phi F+m^\ast\phi^\ast F^\ast
-\frac12m\psi\psi
-\frac12m^\ast\bar\psi\bar\psi.
\label{eq:l8-mass-offshell}
\end{align}
Фермионный билинейный член является майорановской массой для вейлевского поля. Уравнения для \(F\) алгебраические:
\begin{equation}
F^\ast=-m\phi,
\qquad
F=-m^\ast\phi^\ast.
\label{eq:l8-mass-F-equations}
\end{equation}
Подстановка даёт
\begin{equation}
\begin{aligned}
\mathcal L_{\mathrm{on\ shell}}={}&
\partial_\mu\phi^\ast\partial^\mu\phi
+i\bar\psi\bar\sigma^\mu\partial_\mu\psi\\
&-|m|^2|\phi|^2\\
&-\frac12\left(m\psi\psi+m^\ast\bar\psi\bar\psi\right).
\end{aligned}
\label{eq:l8-mass-onshell}
\end{equation}
Следовательно, масса скаляра и масса фермиона обе равны \(|m|\). В использованном на доске диаграммном языке скаляр локально превращается в \(F\) посредством одной вставки \(m\phi F\) и немедленно превращается обратно через сопряжённую вставку. Две вставки дают \(|m|^2|\phi|^2\), то есть в точности скалярный массовый член, найденный при решении уравнения ~\eqref{eq:l8-mass-F-equations}.

Теперь нормировку можно сформулировать точно.\editorialnote{В лекции при перемножении суперполей численные коэффициенты намеренно опускаются. Коэффициенты \(1/2\), \(1/3\) и \(\sqrt2\), использованные здесь, соответствуют стандартной нормировке Wess--Zumino; они делают компонентные массы и константы связи точными, сохраняя при этом каждый качественный шаг вывода на доске.}

\section{Кубический суперпотенциал и взаимодействия}

Теперь добавим
\begin{equation}
W(\Phi)=\frac12m\Phi^2+\frac13g\Phi^3.
\label{eq:l8-superpotential}
\end{equation}
Для общего голоморфного \(W\) \(F\)-компонента имеет вид
\begin{equation}
\left.W(\Phi)\right|_{\theta^2}
=F W'(\phi)-\frac12W''(\phi)\psi\psi.
\label{eq:l8-W-components}
\end{equation}
Здесь
\begin{equation}
W'(\phi)=m\phi+g\phi^2,
\qquad
W''(\phi)=m+2g\phi.
\end{equation}
До исключения \(F\) кубический член содержит
\begin{equation}
g\phi^2F-g\phi\psi\psi+\text{h.c.}
\label{eq:l8-cubic-offshell}
\end{equation}
Второй член — это вершина Юкавы, соединяющая один скаляр с двумя фермионными линиями. Первый соединяет два скаляра со вспомогательной линией.

\begin{figure}[H]
\centering
\begin{tikzpicture}[scale=0.95,>=Latex]
  \coordinate (vone) at (0,0);
  \draw[->,thick] (-2,1.1) -- (vone);
  \draw[->,thick] (vone) -- (-2,-1.1);
  \draw[thick,dashed] (vone) -- (2,0);
  \node[left] at (-2,1.1) {\(\psi\)};
  \node[left] at (-2,-1.1) {\(\psi\)};
  \node[right] at (2,0) {\(\phi\)};
  \node[below=8pt of vone] {\(-g\phi\psi\psi\)};
\end{tikzpicture}
\caption{Наглядная перерисовка юкавского взаимодействия: сходятся две фермионные ножки и одна скалярная, с константой связи, пропорциональной \(g\).}
\end{figure}

Полное уравнение для вспомогательного поля имеет вид
\begin{equation}
F^\ast=-W'(\phi)=-(m\phi+g\phi^2).
\label{eq:l8-general-F-equation}
\end{equation}
Его исключение порождает скалярный потенциал
\begin{equation}
V(\phi)=|W'(\phi)|^2
=|m\phi+g\phi^2|^2.
\label{eq:l8-scalar-potential}
\end{equation}
Помимо массового члена, здесь содержатся кубические скалярные члены, пропорциональные \(m^\ast g\), и квартический член
\begin{equation}
V(\phi)\supset |g|^2|\phi|^4.
\label{eq:l8-quartic}
\end{equation}
В картине со вспомогательной линией две вершины \(\phi^2F\) соединены локальной \(F\)-контракцией. Поскольку уравнение ~\eqref{eq:l8-F-contact} заставляет их положения совпасть, результатом в точности является контактная четырёхскалярная вершина.

\begin{figure}[H]
\centering
\begin{tikzpicture}[scale=0.95]
  \coordinate (vfour) at (0,0);
  \foreach \angle in {35,145,215,325}
    \draw[thick,dashed] (vfour) -- (\angle:1.9);
  \fill (vfour) circle (1.5pt);
  \node[below=8pt of vfour] {\(-|g|^2|\phi|^4\)};
\end{tikzpicture}
\caption{После исключения вспомогательного поля две локальные \(F\)-вершины превращаются в одно четырёхскалярное контактное взаимодействие.}
\end{figure}

Суперсимметрия сделала больше, чем просто ввела бозон и фермион в одну и ту же теорию. Она заставила их массы совпасть, связала юкавскую константу связи с \(g\) и зафиксировала константу четырёхскалярного взаимодействия равной \(|g|^2\). Эти соотношения можно было бы проверить по массам, амплитудам рассеяния и распадам; это не необязательные численные совпадения.

\section{Парные петлевые эффекты}

Эти соотношения становятся особенно полезными, когда взаимодействия вносят поправки в массы частиц. В собственную энергию скаляра входят бозонная петля от квартической вершины и фермионная петля от двух юкавских вершин. Каждая из них содержит две степени одной и той же константы связи. Замкнутая фермионная петля имеет относительный знак минус, требуемый статистикой Ферми.

\begin{figure}[H]
\centering
\begin{tikzpicture}[>=Latex,scale=0.9]
  \begin{scope}[xshift=-3.3cm]
    \draw[thick,dashed] (-2,0)--(0,0)--(2,0);
    \draw[thick,dashed] (0,0) .. controls (-0.9,1.7) and (0.9,1.7) .. (0,0);
    \fill (0,0) circle (1.4pt);
    \node at (0,-0.55) {bosonic loop};
  \end{scope}
  \begin{scope}[xshift=3.3cm]
    \draw[thick,dashed] (-2,0)--(-0.8,0);
    \draw[thick,dashed] (0.8,0)--(2,0);
    \draw[->,thick] (-0.8,0) arc[start angle=180,end angle=0,x radius=0.8,y radius=0.8];
    \draw[->,thick] (0.8,0) arc[start angle=0,end angle=-180,x radius=0.8,y radius=0.8];
    \fill (-0.8,0) circle (1.4pt);
    \fill (0.8,0) circle (1.4pt);
    \node at (0,-1.35) {fermionic loop};
  \end{scope}
\end{tikzpicture}
\caption{Рассматриваемые в лекции парные механизмы вклада в скалярную собственную энергию. Суперсимметрия фиксирует относительные константы связи, а замкнутая фермионная петля даёт противоположный знак.}
\end{figure}

На одном петлевом уровне ведущие ультрафиолетово-чувствительные части сокращаются в сумме полных бозонного и фермионного вкладов. Та же симметрия организует высшие порядки посредством тождеств Уорда и свойств ненормировки. Это утверждение относится к суммам, связанным симметрией, а не к буквальному взаимно-однозначному сокращению каждой диаграммы; могут сохраняться перенормировка волновой функции, конечные эффекты и другие поправки, сохраняющие симметрию. В этом и состоит практическая сила построения теории в суперпространстве: равенство масс и соотношения между константами связи обеспечиваются ещё до вычисления какого бы то ни было отдельного интеграла Фейнмана.\editorialnote{На доске это сокращение сформулировано более широко; в основном тексте дана стандартная оговорка.}

\section{Что дальше}

Построенная здесь модель — простейшая взаимодействующая киральная теория, обычно называемая моделью Wess--Zumino. Это пример, а не наиболее общая суперсимметричная полевая теория, но общая схема сохраняется: поля представляются супермультиплетами, записываются инвариантные интегралы по суперпространству, а разложение по переменным Грассмана навязывает соотношения между компонентами.

Точная суперсимметрия не видна в наблюдаемом спектре частиц, поэтому следующая проблема состоит в том, как можно нарушить эту симметрию, сохранив при этом достаточно контроля над ультрафиолетовым поведением, чтобы это оставалось полезным. Другое обещанное направление — великое объединение: как
\begin{equation}
SU(3)\times SU(2)\times U(1)
\end{equation}
может вписываться в более крупную группу, такую как \(SU(5)\), и как суперсимметричный набор частиц меняет эту картину объединения. Именно к этим вопросам обращаются оставшиеся лекции.
